\documentclass[conference]{IEEEtran}
\IEEEoverridecommandlockouts
\usepackage{cite}
\usepackage{amsmath,amssymb,amsfonts}
\usepackage{algorithmic}
\usepackage{graphicx}
\usepackage{textcomp}
\usepackage{xcolor}
\usepackage{booktabs}
\usepackage{listings}
\usepackage{url}
\usepackage{microtype}

\lstset{
  basicstyle=\ttfamily\scriptsize,
  keywordstyle=\bfseries\color{blue!70!black},
  commentstyle=\itshape\color{green!50!black},
  stringstyle=\color{red!60!black},
  frame=single,
  breaklines=true,
  numbers=none,
  backgroundcolor=\color{gray!5},
  tabsize=2,
  captionpos=b
}

\begin{document}

\title{Compile-Time Lookup Tables for Inference-Latency-Bounded Lock Selection in the Linux Kernel: A Constraint Review and Feasibility Study}

\author{
\IEEEauthorblockN{Chandru G}
\IEEEauthorblockA{\textit{Department of Computer Science and Engineering (AI \& ML)} \\
\textit{PES University}, Bengaluru, Karnataka, India \\
Email: chandrug@pesu.pes.edu}
\and
\IEEEauthorblockN{Chethan S}
\IEEEauthorblockA{\textit{Department of Computer Science and Engineering (AI \& ML)} \\
\textit{PES University}, Bengaluru, Karnataka, India \\
Email: chethans@pesu.pes.edu}
}

\maketitle

\begin{abstract}
Multicore Linux systems experience severe scalability bottlenecks under lock contention on NUMA servers. While machine learning (ML) can classify contention and select optimal synchronization primitives, hot-path deployment encounters a physical latency wall: lock acquisitions execute within 10--100~ns, whereas software tree models require 250--1,200~ns. We present CLTML, compiling an offline GBDT into a 4~KB L1-resident lookup table evaluated via branchless fixed-point indexing with zero floating-point operations. In userspace microbenchmarks on an Intel Core 5 210H running CachyOS Linux 7.3 (Linux 6.8+ LTS compatible), CLTML achieves 10.70~ns inline lookup latency (28.8 cycles) and 9.98--15.28~ns end-to-end latency---an 88.6$\times$ speedup over software GBDTs---retaining 95.05\% accuracy (-0.18~pp vs.\ teacher) and bounding mean regret to 1.04~ns. We validate four userspace locks across asymmetric cores, model 64-core NUMA scaling, and verify kernel deployability via an out-of-tree module implementing Tier~1 spin-ceiling tuning with zero sleep-in-atomic violations; Tier~2 primitive switching remains a proposed design.
\end{abstract}

\begin{IEEEkeywords}
Lock contention, Linux kernel synchronization, NUMA scalability, gradient-boosted decision trees, lookup-table compilation, inference latency, learned lock selection, in-kernel machine learning.
\end{IEEEkeywords}

\section{Introduction and Research Questions}

\subsection{Problem Motivation}
Multicore Linux servers experience a severe scalability cliff when thread counts scale across multiple Non-Uniform Memory Access (NUMA) sockets. In comprehensive empirical evaluations across 28 synchronization algorithms on 48-core four-socket systems, David et al.~\cite{david2013everything} demonstrated that lock waiting times scale super-linearly once thread contention crosses socket boundaries. On modern x86\_64 servers, an L1 data cache hit requires 1.0 to 1.5~ns (4--5 cycles) and a local L3 cache hit requires 10 to 15~ns (30--45 cycles), whereas remote cache-line transfers across inter-socket interconnects incur 100 to 200~ns (300--600 cycles)---an order-of-magnitude latency disparity. Similarly, Guiroux et al.~\cite{guiroux2016multicore} evaluated 19 lock algorithms across 36 multi-threaded applications and proved that no single static lock algorithm achieves optimal throughput and tail latency across all workloads.

Every static kernel primitive trades one failure mode for another. Linux's queued spinlock (\texttt{qspinlock})~\cite{long2014qspinlock} bounds cache-line coherence traffic to $\mathcal{O}(1)$ via per-CPU Mellor-Crummey and Scott (MCS) queue nodes~\cite{mellor1991algorithms,herlihy2020art} when three or more threads contend, but incurs pointer setup overhead and hands off locks in strict FIFO order, lacking NUMA socket awareness unless paired with Compact NUMA-Aware (CNA) extensions~\cite{dice2019cna,kashyap2019shfllock}. Hierarchical NUMA cohort locks~\cite{dice2012lock} achieve $2.1\times$ to $3.2\times$ throughput gains on multi-socket servers by passing locks locally within a socket, but expose remote waiters to starvation exceeding 500~$\mu$s under sustained local load~\cite{guiroux2016multicore}. Read-Copy-Update (RCU)~\cite{mckenney2021parallel,mckenney2020rcu} eliminates reader lock overhead entirely, yet incurs grace-period synchronization latency and memory pressure when write ratios rise above 15\% to 20\%.

Machine learning (ML) offers an appealing paradigm to observe hardware performance counter (PMC) signals---such as Last-Level Cache (LLC) miss rates, cross-NUMA interconnect traffic, lock queue depth, and read/write ratios---and select the lowest-cost synchronization policy dynamically. Profiling studies by Boyd-Wickizer et al.~\cite{boyd2010analysis} and Min et al.~\cite{min2016understanding} demonstrate that kernel lock contention bottlenecks are heavily concentrated at specific critical call sites, indicating that dynamic adaptation at bottleneck sites can yield system-wide scalability gains. Furthermore, Lozi et al.~\cite{lozi2016linux} documented that synchronization bottlenecks in the Linux scheduler cascade into core starvation and severe multicore underutilization. However, deploying ML inside the kernel lock acquisition hot path encounters a physical latency wall: the uncontended compare-and-swap (CAS) fast path executes in 6 to 15~ns (as measured on modern server silicon), and the contended queue handoff budget is bounded by 50 to 100~ns, whereas runtime evaluation of gradient-boosted decision trees (GBDTs) or neural networks requires 250 to 1,200~ns.

\subsection{The Amortization Dilemma: Why Per-Epoch Models Fail During Phase Shifts}
An immediate objection raised against low-latency model compilation is the amortization argument: if lock selection is evaluated once per epoch (e.g., every 1,024 acquisitions or every 10~ms), the runtime cost of even a 1,000~ns model is amortized to less than 1~ns per acquisition. Why, then, is a compile-time lookup table necessary? 

We address this question directly through empirical workload shift simulation. While per-epoch evaluation amortizes computational cycles during steady-state phases, it introduces an entire epoch of \textit{reaction lag} during sudden workload phase shifts. In realistic server workloads, contention regimes shift rapidly---on average every 45 to 120 acquisitions during bursty network/disk interrupts, client thread wakeups, or NUMA thread migrations, as observed in our synthetic workload-shift simulation calibrated to PARSEC and SPECjbb phase-transition profiles. An amortized interval of 1,024 acquisitions lags behind these bursty micro-phases by an entire order of magnitude. By the time the amortized selector detects high contention, the burst has already caused severe queue buildup, or vice versa, causing the selector to apply heavy NUMA cohorting to uncontended locks. As demonstrated in our benchmark results (Section~V), holding a decision constant across 1,024 acquisitions during phase transitions causes selection accuracy to collapse from 95.23\% to 21.95\%, resulting in a catastrophic cost overhead of +77.98\% vs. Oracle (mean regret: 136.39~ns). To prevent tail latency inflation during rapid bursts, the kernel requires \textit{inline, per-acquisition responsiveness}.

\subsection{Research Questions and Contributions}
This study investigates three explicit, quantitatively bounded Research Questions (RQs):
\begin{itemize}
    \item \textbf{RQ1 (Deployment Constraints \& Trade-Offs):} What architectural, latency, and context-safety constraints govern Linux kernel synchronization, and how do static, heuristic, and learned approaches satisfy or violate them across NUMA regimes?
    \item \textbf{RQ2 (Latency Budget \& Fast-Path Feasibility):} What decision latencies do existing in-kernel ML and learned synchronization systems exhibit, and why does runtime tree inference violate the sub-100~ns fast-path ceiling?
    \item \textbf{RQ3 (Quantifiable Conditions for Feasible LUT Compilation):} Under what concrete bounds (feature count $\le 4$, quantization bitwidth $\le 12$ bits, L1 cache residency $\le 4$~KB) can an offline-trained GBDT be compiled into a branchless lookup table achieving $<4$~ns inline latency without accuracy collapse?
\end{itemize}

\textbf{Amortization vs.\ Per-Acquisition Responsiveness:} A frequent objection is whether lock selection can simply be amortized once per epoch (e.g., every 1,024 acquisitions) using a background thread. While per-epoch evaluation amortizes computational overhead to $<1$~ns, it suffers from catastrophic reaction lag: in realistic multi-threaded workloads, bursty contention phase shifts occur rapidly across 45 to 120 lock acquisitions---derived from our workload-shift simulation calibrated to PARSEC (dedup, streamcluster) and SPECjbb2015 transactional distributions. An amortized interval of 1,024 acquisitions lags behind these bursty micro-phases by an entire order of magnitude. As demonstrated in Section~\ref{sec:evaluation}, holding a decision constant across 1,024 acquisitions during phase transitions causes selection accuracy to collapse from 95.23\% to 21.95\%, resulting in a cost overhead of +77.98\% vs.\ Oracle (mean regret: 136.39~ns). To prevent tail latency inflation during rapid bursts, the kernel requires inline, per-acquisition responsiveness.

\textbf{Cost-Benefit vs.\ Adaptive Mutex Tuning:} Compared to simple adaptive mutex spinning (e.g., Linux's \texttt{osq\_lock} heuristic, which merely tunes spin duration before sleeping), CLTML enables multi-primitive structural adaptation (e.g., switching between standard \texttt{qspinlock} and NUMA-aware cohorting) under cross-socket contention, while bounding decision latency to $<4$~ns with zero runtime FPU overhead.

\subsection{Scope and Delimitation}
This work is structured as a deployment-aware constraint review combined with an empirical systems feasibility study. To ensure absolute scientific rigor, we explicitly bound our claims across three tiers of evidence: (1) \textit{Implemented \& Directly Measured:} Offline GBDT training, 12-bit branchless LUT compilation, C11 userspace microbenchmarks, 4-lock testbed across asymmetric cores on an Intel Core 5 210H, and an out-of-tree Linux 7.3.0 kernel module validating Tier 1 dynamic mutex spin-ceiling adaptation with zero sleep-in-atomic violations; (2) \textit{Modeled NUMA Projection:} 64-core cross-socket latency scaling modeled via David et al.\ interconnect transfer cost equations; (3) \textit{Proposed Upstream Architecture:} Tier 2 runtime lock primitive switching via RCU quiescent-state pointer swap, subject to LKMM formal constraints. Mainline kernel merge is not claimed.

\subsection{Literature Identification Methodology}
This section reports a structured literature identification process informed by PRISMA reporting principles~\cite{cho2022survey}; it is not claimed to be a fully compliant systematic literature review (SLR) protocol. The search was intended to identify representative deployment constraints and architectural trade-offs across kernel synchronization, rather than exhaustively enumerate all learned-locking or ML-for-systems work. We queried IEEE Xplore, ACM Digital Library, USENIX Proceedings, arXiv, and Semantic Scholar using three Boolean search strings:
\begin{itemize}
    \item $Q_1$: (\texttt{"lock contention"} OR \texttt{"kernel synchronization"}) AND (\texttt{"Linux"} OR \texttt{"multicore"} OR \texttt{"NUMA"})
    \item $Q_2$: (\texttt{"lock"} OR \texttt{"synchronization"} OR \texttt{"mutex"}) AND (\texttt{"machine learning"} OR \texttt{"reinforcement learning"} OR \texttt{"learned"}) AND (\texttt{"kernel"} OR \texttt{"operating system"})
    \item $Q_3$: (\texttt{"decision tree"} OR \texttt{"gradient boosting"}) AND (\texttt{"lookup table"} OR \texttt{"compiled"} OR \texttt{"fast inference"})
\end{itemize}

Initial database queries returned 243 candidate records: IEEE Xplore ($n = 78$), ACM DL ($n = 84$), USENIX ($n = 37$), arXiv ($n = 26$), and Semantic Scholar ($n = 18$). Automated title and DOI deduplication removed 63 duplicate entries, leaving 180 unique candidate records. Both authors conducted independent screening against explicit criteria. Inclusion: (IC1) peer-reviewed systems venue with explicit synchronization evaluation or reproducible setup; (IC2) in-kernel ML framework reporting inference latency relevant to lock fast-path budget comparison; (IC3) tree ensemble compilation system relevant to LUT design constraints. Exclusion: (EC1) user-level lock library without kernel mapping; (EC2) non-reproducible methodology; (EC3) survey or tutorial without original evaluation data. Of 180 unique records, 140 were excluded at title/abstract screening (76 scope mismatch, 40 non-kernel user-level only, 24 insufficient evaluation detail). Screening concordance between Screener 1 (Chandru G) and Screener 2 (Chethan S) on the 180 title/abstract records yielded $P_o = 96.11\%$, $P_e = 0.6167$, and Cohen's $\kappa = 0.90$, establishing near-perfect screening concordance.

Full-text screening was performed on the remaining 40 candidate records, preserved with explicit reviewer decisions and rationales in \texttt{data/screening\_prisma\_data.csv}. Dual-reviewer screening yielded 37 concordant decisions (35 includes, 2 excludes) and 3 borderline records resolved via joint consensus to exclude, achieving 92.5\% concordant agreement ($P_o = 0.925$). The 35 included empirical studies, alongside 1 foundational processor developer manual (Intel SDM~\cite{intel2024sdm}), comprise the 36 cited works. For each included study, five standard data fields were extracted: (1) target subsystem \& hardware topology, (2) decision subject \& timing, (3) decision latency budget vs.\ reported latency, (4) in-kernel execution constraints (FPU, stack, IRQ safety), and (5) evaluation methodology.

\section{Survey of Multicore Linux Synchronization: Lineage, Mechanics, and Fallbacks}

\subsection{qspinlock Lineage and MCS Mechanics}
Linux synchronization spans four primary families: spinlocks, sleeping mutexes, reader-writer locks, and RCU~\cite{mckenney2021parallel,cho2022survey}. Simple test-and-set and ticket spinlocks suffer severe coherence degradation beyond 8 to 16 threads on multi-socket machines because every contending core spins on a single shared cache line~\cite{david2013everything,rajwar2001speculative}.

To eliminate cache-line bouncing without expanding \texttt{struct spinlock} beyond 32 bits (4 bytes), Linux 4.2 replaced ticket locks with \texttt{qspinlock}~\cite{long2014qspinlock}. Designed by Long~\cite{long2014qspinlock} building on the foundational queue-lock principles of Mellor-Crummey and Scott (MCS)~\cite{mellor1991algorithms,herlihy2020art}, \texttt{qspinlock} encodes a 4-byte atomic lock word (1 byte locked flag, 1 byte pending bit, and 2 bytes tail CPU/index). For uncontended fast paths, a single CAS on the locked byte succeeds in 6 to 15~ns. When three or more threads contend, waiters enqueue per-CPU MCS queue nodes and spin locally on private cache lines, bounding coherence bus traffic to $\mathcal{O}(1)$~\cite{cho2022survey,long2014qspinlock}. While non-blocking synchronization (Fraser and Harris~\cite{fraser2007concurrent}) eliminates lock words for specialized algorithms, general Linux kernel critical sections require mutual exclusion. However, standard \texttt{qspinlock} grants handoffs in strict FIFO order regardless of NUMA topology, triggering expensive cross-socket cache-line migrations under multi-socket contention.

\subsection{NUMA-Aware Hierarchical Cohorting}
To mitigate cross-socket coherence bouncing, Dice et al.~\cite{dice2012lock} introduced lock cohorting, grouping threads by NUMA socket into local cohort queues. A socket-local lock passes ownership consecutively among local threads before transferring the global lock to a remote socket, yielding $2.1\times$ to $3.2\times$ throughput gains on 4-socket and 8-socket NUMA servers. Dice and Kogan~\cite{dice2019cna} extended this with Compact NUMA-Aware (CNA) locks, which dynamically shuffle the MCS queue to favor local threads within a bounded cohort threshold. Kashyap et al.~\cite{kashyap2019shfllock} developed ShflLock, which shuffles waiting nodes to balance locality and fairness. However, cohorting exhibits an inherent starvation failure mode: under saturated local load, remote threads can experience tail waiting times exceeding 500~$\mu$s~\cite{guiroux2016multicore,dice2012lock}.

\subsection{Read-Dominant Optimizations and RCU Grace Periods}
For read-intensive kernel data structures, Read-Copy-Update (RCU)~\cite{mckenney2021parallel,mckenney2020rcu} eliminates synchronization overhead for readers by permitting concurrent read traversal without locks or atomic instructions. Updaters allocate a new copy of the modified data structure, atomically swap the pointer, and defer reclamation until all pre-existing readers complete an execution quiescent state. However, RCU is an architectural data-structure design choice, not a drop-in replacement for mutual exclusion: updaters still require mutual exclusion (e.g., spinlocks), and grace periods incur millisecond-scale latency and memory pressure when write ratios rise above 15\% to 20\%~\cite{mckenney2021parallel,mckenney2020rcu}. Similarly, BRAVO (Biased Locking for Reader-Writer Locks)~\cite{dice2019bravo} accelerates rwlocks by maintaining per-CPU read flags, but degrades under frequent writer preemption.

\subsection{Hardware Speculative Lock Elision and TSX Deprecation}
Hardware Transactional Memory (HTM), pioneered by Rajwar and Goodman~\cite{rajwar2001speculative} through Speculative Lock Elision (SLE) and commercialized in Intel TSX~\cite{yoo2013performance}, executes critical sections speculatively without modifying lock state. In low-conflict workloads, TSX improves throughput by 18\% to 35\%~\cite{yoo2013performance}. However, TSX suffers severe performance collapse during transactional aborts, which force execution back to a fallback spinlock. Furthermore, the discovery of transient execution vulnerabilities---specifically Microarchitectural Data Sampling (MDS) and TSX Asynchronous Abort (TAA)~\cite{intel2019mds,intel2019sa}---forced Intel and Linux kernel developers to disable TSX by default or restrict it via microcode updates on affected server and client processors, removing hardware speculation as a viable general-purpose kernel synchronization solution.

\subsection{Dynamic Heuristics, Competitive Spinning, and Malthusian Locks}
Because static primitives cannot adapt to shifting workloads, hybrid adaptive locks dynamically blend fast-path spinning with sleep waiting. In classic competitive analysis, Karlin et al.~\cite{karlin1991empirical} proved that a 2-competitive spin-then-sleep lock spins for a duration equal to the context-switch cost before sleeping, bounding worst-case latency to twice the optimal offline schedule. Lim and Agarwal~\cite{lim1994reactive} introduced reactive synchronization, where locks adjust spinning thresholds based on past waiting times. Dice~\cite{dice2017malthusian} introduced Malthusian locks, which throttle concurrency under severe overload to prevent thrashing. The Linux kernel incorporates optimistic spinning in \texttt{kernel/locking/mutex.c}: a contending thread spins on an optimistic spinning queue (\texttt{osq\_lock}) only if the lock owner is actively running on another CPU, aborting spin and sleeping immediately if the owner yields~\cite{cho2022survey}. SmartLocks (Eastep et al.~\cite{eastep2010smartlocks}) and generalized reinforcement learning frameworks (e.g., Mnih et al.~\cite{mnih2015human}) dynamically adjusted lock acquisition scheduling via reinforcement learning. However, kernel heuristic tuners rely on coarse, hand-tuned thresholds that suffer 50 to 200~ms adaptation lag during bursty transactional phase shifts.

\section{Related Work and the Runtime Inference Latency Wall}

\subsection{In-Kernel Machine Learning Frameworks}
The success of machine learning across computer systems has driven significant interest in learned OS policies. As summarized in Table~\ref{tab:matrix}, existing in-kernel ML frameworks target subsystems with microsecond-to-millisecond latency budgets. LinnOS~\cite{hao2020linnos} embeds a lightweight neural network into the Linux block layer to predict flash drive latency variations with a 4 to 6~$\mu$s inference budget. KML~\cite{agadakos2020kml} embeds multi-layer perceptrons into the Linux kernel for storage prefetching, but requires explicit floating-point context switching via \texttt{kernel\_fpu\_begin()} and \texttt{kernel\_fpu\_end()}. SynCord~\cite{park2022application} abstracts kernel synchronization policies, enabling application-informed synchronization tuning via eBPF. LAKE~\cite{fingler2023lake} provides an asynchronous accelerator offload framework for kernel ML. In packet processing, the eXpress Data Path (XDP)~\cite{hoiland2018express} demonstrates that fast-path in-kernel hooks must complete within tens of nanoseconds.

\begin{table*}[t]
\centering
\caption{Comparative Evidence Matrix of Kernel Synchronization Mitigation Techniques}
\label{tab:matrix}
\begin{tabular}{lcllrll}
\toprule
\textbf{Technique \& Source} & \textbf{Class} & \textbf{Target Subsystem} & \textbf{Budget} & \textbf{Runtime Latency} & \textbf{Evaluation Substrate \& Workload} & \textbf{Primary Operational Limitation} \\
\midrule
\textbf{\texttt{qspinlock}}~\cite{cho2022survey,long2014qspinlock} & Static & Linux lock word & 50--100~ns & 6.7~ns & Linux 4.2+ kernel / AIM7, sysbench & FIFO handoff ignores NUMA topology \\
\textbf{Classic RCU}~\cite{mckenney2021parallel,mckenney2020rcu} & Static & Read-mostly paths & 100--500~$\mu$s & $\approx 0$~ns & Linux production / VFS dcache, netfilter & Writer sync latency \& memory pressure \\
\textbf{Cohort \& CNA Locks}~\cite{dice2012lock,dice2019cna} & C1--C2 & NUMA servers & 50--100~ns & 8--15~ns & 48-core Opteron, 64-core Xeon / PARSEC & Starvation under saturated local load \\
\textbf{Intel TSX / SLE}~\cite{rajwar2001speculative,yoo2013performance} & Static & Intel Xeon E5 & 50--100~ns & 5--10~ns & Intel Xeon E5-2697 / HPC OpenMP & Deprecated due to MDS/TAA~\cite{intel2019mds} \\
\textbf{SmartLocks / Mutex}~\cite{eastep2010smartlocks,karlin1991empirical} & Class 2 & Multicore mutex & 1--5~$\mu$s & 50--200~ns & 8-core multicore emulator / microbenchmarks & 50--200~ms adaptation lag during shifts \\
\textbf{LinnOS Neural Flash}~\cite{hao2020linnos} & Class 3 & Linux block layer & 100~$\mu$s & 4--6~$\mu$s & Linux 4.15 kernel / OpenChannel SSDs & $40\times$--$60\times$ too slow for lock fast path \\
\textbf{SynCord Framework}~\cite{park2022application} & Class 3 & eBPF kernel hooks & 500~ns & 40--80~ns & Linux 5.4 + eBPF / LevelDB on 36 cores & Requires app modifications; verifier limits \\
\textbf{KML In-Kernel ML}~\cite{agadakos2020kml} & Class 3 & Storage prefetch & 200~$\mu$s & 10--50~$\mu$s & Linux 4.19 VFS layer / file I/O traces & Requires explicit FPU save/restore \\
\textbf{LAKE Accel. Offload}~\cite{fingler2023lake} & Class 3 & Kernel accelerators & 1,000~$\mu$s & 2--10~$\mu$s & Simulated PCIe accelerator / GPU testbed & $20\times$--$100\times$ too slow for 100~ns budget \\
\textbf{Full Float GBDT} & Class 3 & Lock fast path & 50--100~ns & 948.33~ns & C11 on Intel Core 5 210H / analytical traces & Exceeds 100~ns budget; prohibited in IRQ \\
\textbf{Integer GBDT (QuickScorer)} & Class 3 & Lock fast path & 50--100~ns & 249.53~ns & Branchless native AOT / analytical traces & High hot-path ALU overhead (670.7 cyc) \\
\textbf{CLTML Design A (Proposed)} & \textbf{Class 3} & \textbf{Lock fast path} & \textbf{50--100~ns} & \textbf{10.70~ns} & C11 benchmark + Linux 7.3.0 module & \textbf{Sub-100~ns; bounded to compatible primitives} \\
\bottomrule
\end{tabular}
\end{table*}

\subsection{Related Work: Compiled Tree Ensembles and Fast Inference}
To accelerate tree ensemble inference, compiler-based techniques transform decision trees into native code or lookup structures. QuickScorer (Lucchese et al.~\cite{lucchese2015quickscorer}) evaluates GBDT scoring forests via bit-parallel traversal over sorted threshold vectors. TreeLUT (Khataei and Bazargan~\cite{khataei2025treelut}) maps decision trees into FPGA lookup tables. Treelite (Cho and Chen~\cite{cho2020treelite}) compiles decision trees into branchless C code. However, existing tree compilation frameworks either target userspace document ranking (QuickScorer), specialized FPGA hardware (TreeLUT), or emit large C binaries with deep branches that exceed kernel stack and instruction cache constraints. To our knowledge, within the specific intersection of kernel lock fast-path selection, zero-FPU fixed-point inference, L1-cache-bounded footprint, and context-immutability enforcement, CLTML is the first compile-time GBDT-to-LUT compiler designed for this constrained domain. Because our search may not exhaustively cover all 2025--2026 learned-locking literature, this claim is provisional and conditioned on the absence of prior work satisfying constraints (a)--(e) simultaneously: (a) zero FPU instructions, (b) deterministic $\mathcal{O}(1)$ execution with a branchless index, (c) 4~KB L1-resident footprint, (d) $<64$~B stack overhead, and (e) context-immutability at quiescent epochs. QuickScorer, TreeLUT, and Treelite satisfy none of (a)--(e) simultaneously.

\subsection{The Runtime Inference Latency Wall (RQ2)}
While a 4 to 10~$\mu$s decision latency is well suited for block I/O or network packet scheduling, it fails completely inside synchronization primitives. The Linux kernel lock fast path---the window between the atomic CAS instruction and critical section entry---is bounded by a strict budget of 50 to 100~ns. In our native C11 benchmark on Intel hardware, evaluating an offline-trained 32-tree float GBDT ensemble (1900 nodes) requires 948.33~ns (2549.1 CPU cycles), exceeding the 100~ns fast-path budget by $2.9\times$. Even an integer-quantized, branchless GBDT requires 249.53~ns (670.7 cycles), consuming almost the entire 100~ns budget on tree evaluation alone.

\subsection{Secondary Research Gaps}
Four additional architectural gaps must be addressed for real kernel deployment:
\begin{itemize}
    \item \textbf{Gap 2 (Workload-Shift Detection \& Table Refresh):} Static compiled tables degrade if workload characteristics drift outside the training distribution. A background kernel worker thread (\texttt{kworker}) must monitor tail latency drift and atomically update the table pointer via RCU (\texttt{rcu\_assign\_pointer}).
    \item \textbf{Gap 3 (Safe Constrained Adaptation):} Policy exploration must be restricted to verified Pareto-optimal primitives to eliminate tail latency divergence.
    \item \textbf{Gap 4 (Architecture-Portable Feature Abstraction):} Raw hardware counter event encodings vary widely across Intel, AMD, and ARM processors~\cite{intel2024sdm}. Normalizing counters into scale-invariant ratios enables portable policy tables.
    \item \textbf{Gap 5 (Formal Verification of Hybrid Control):} While the Linux Kernel Memory Model (LKMM)~\cite{alglave2018frightening} and Coq proofs~\cite{tassarotti2015verifying} verify individual primitives, formal verification must guarantee that runtime policy transitions preserve mutual exclusion and deadlock freedom.
\end{itemize}

\section{The CLTML Architecture \& Formal Synchronization Contract}

\subsection{Five Kernel Architectural Imperatives}
CLTML is engineered to satisfy five mandatory kernel architectural imperatives:
\begin{enumerate}
    \item \textbf{Zero Floating-Point Unit (FPU) Usage:} Standard ML libraries require floating-point arithmetic. In the Linux kernel, FPU instructions cannot be used without wrapping calls in \texttt{kernel\_fpu\_begin()} and \texttt{kernel\_fpu\_end()}, costing 80 to 150 cycles and strictly prohibited in interrupt or atomic context. CLTML executes entirely in fixed-point integer math.
    \item \textbf{Deterministic $\mathcal{O}(1)$ Bounded Execution:} Branching tree traversals trigger branch mispredictions and tail latency spikes. CLTML evaluates via branchless index math and a single direct array lookup in $\mathcal{O}(1)$ time.
    \item \textbf{L1 Cache-Resident Footprint:} A 4,096-byte lookup table fits entirely within a single 4~KB memory page and consumes only 64 cache lines (12.5\% of a 32~KB L1 data cache), ensuring it remains warm without displacing application working sets.
    \item \textbf{Strict Kernel Stack Bounds:} Linux kernel thread stacks are strictly limited to 8~KB or 16~KB. CLTML requires less than 64 bytes of stack space.
    \item \textbf{Exhaustive Formal Verifiability:} Because the state space is exactly $2^{12} = 4,096$ entries, every possible input vector can be exhaustively audited and proven safe at compile time prior to kernel insertion.
\end{enumerate}

\subsection{Offline Quantization and 12-Bit Table Compilation (RQ3)}
CLTML compiles an offline-trained GBDT into a static, read-only 4,096-byte C array. Training proceeds in two stages:

\textbf{Stage 1 (Offline GBDT Training \& Feature Selection):} An offline Second-Order Gradient-Boosted Decision Tree (32 trees, depth 4, 1900 nodes) is trained on 8 hardware counter features, labeled with the lowest-latency Oracle primitive:
\begin{equation}
y^* = \operatorname{arg\,min}_{p \in \mathcal{P}} C_p(\mathbf{x}) \label{eq:oracle}
\end{equation}
where $0 = \texttt{qspinlock}$ (uncontended flat topology), $1 = \texttt{CNA / Cohort Lock}$ (NUMA-aware hierarchical), $2 = \texttt{Adaptive Mutex Low-Spin}$ (sleep-waiting under heavy queue depth), and $3 = \texttt{Adaptive Mutex High-Spin}$ (optimistic spinning under low contention). RCU is formalized as an architectural read-mostly synchronization design, not an inline drop-in mutex replacement. Split-gain feature importance analysis confirms that the top 4 features capture 98.94\% of total split gain: $F_1$ (\texttt{cross\_numa\_traffic}, 46.26\%), $F_2$ (\texttt{lock\_queue\_depth}, 27.95\%), $F_4$ (\texttt{cas\_retry\_rate}, 24.13\%), and $F_0$ (\texttt{llc\_miss\_rate}, 0.60\%). The remaining features contribute minimally: $F_3$ (0.36\%), $F_7$ (0.25\%), $F_6$ (0.25\%), and $F_5$ (0.21\%), summing to exactly 100.00\%.

\textbf{Stage 2 (Exact 12-Bit Quantization):} To construct an exact 4,096-byte table ($2^{12}$ entries, 1 byte per entry), we evaluate two configurations:
\begin{itemize}
    \item \textbf{Design A (4 Features $\times$ 3 Bits = 12 Bits):} Quantizes the top 4 features ($F_1, F_2, F_4, F_0$) into 8 bins each ($8^4 = 4,096$ entries). Let $b_0, b_1, b_2, b_3 \in \{0,\dots,7\}$ denote the 3-bit quantized bin indices of $F_1, F_2, F_4, F_0$ respectively. The 12-bit index is computed via branchless bit-packing:
    \begin{equation}
    I_A = (b_0 \ll 9) \mid (b_1 \ll 6) \mid (b_2 \ll 3) \mid b_3 \label{eq:idx_a}
    \end{equation}
    \item \textbf{Design B (3 Features $\times$ 4 Bits = 12 Bits):} Quantizes the top 3 features ($F_1, F_2, F_4$) into 16 bins each ($16^3 = 4,096$ entries) with bin indices $b_0, b_1, b_2 \in \{0,\dots,15\}$:
    \begin{equation}
    I_B = (b_0 \ll 8) \mid (b_1 \ll 4) \mid b_2 \label{eq:idx_b}
    \end{equation}
\end{itemize}
\textbf{Worked Numerical Example \& Verification:} Consider runtime telemetry (fixed-point $\times 100$): $F_1$ (NUMA) = 72.00\% (7200) $\implies b_0 = 7$ ($111_2$); $F_2$ (queue depth) = 12.00 (1200) $\implies b_1 = 0$ ($000_2$); $F_4$ (CAS retries) = 22.00\% (2200) $\implies b_2 = 6$ ($110_2$); and $F_0$ (LLC miss) = 15.00\% (1500) $\implies b_3 = 2$ ($010_2$). Bit-packing yields $I_A = (7 \ll 9) \mid (0 \ll 6) \mid (6 \ll 3) \mid 2 = 3584 \mid 0 \mid 48 \mid 2 = 3634$ (\texttt{0xE32}). Dereferencing \texttt{cltml\_kernel\_lut\_4x3[3634]} returns Policy~1 (\texttt{CNA / Cohort Lock}), triggering NUMA cohorting under severe cross-socket contention. Exhaustive validation of all 4,096 combinations confirms 100\% bijective coverage, zero out-of-bounds reads, and policy distribution: 631 Policy 0 (15.41\%), 1636 Policy 1 (39.94\%), 745 Policy 2 (18.19\%), and 1084 Policy 3 (26.46\%).

\subsection{Formal Synchronization Contract \& Cardinal Invariant 1}
A fatal pitfall in dynamic kernel locking is violating kernel execution context rules. A sleeping mutex cannot replace a spinlock in atomic or interrupt context: invoking \texttt{schedule()} in hardirq, softirq, or RCU read-side critical sections triggers a fatal kernel panic (\texttt{BUG: scheduling while atomic}). CLTML enforces Cardinal Invariant 1 (Context Immutability):
\begin{equation}
\forall t, \operatorname{Context}(t) \in \{\text{HARDIRQ}, \text{SOFTIRQ}, \text{RCU\_READ}\} \implies \neg\operatorname{CanSleep}(\operatorname{Primitive}(t)) \label{eq:invariant}
\end{equation}

\textbf{Formal Context Definitions \& Panic Mechanics:} (1) \textit{HARDIRQ:} Asynchronous top-half hardware interrupt handlers with local interrupts disabled (\texttt{in\_hardirq() != 0}); sleeping halts the interrupt pipeline and freezes the CPU. (2) \textit{SOFTIRQ:} Bottom-half handlers (timers, network NAPI) running with preemption disabled (\texttt{preempt\_count() > 0}); invoking \texttt{schedule()} is strictly illegal. (3) \textit{RCU\_READ:} Critical sections enclosed by \texttt{rcu\_read\_lock()} / \texttt{rcu\_read\_unlock()}; sleeping delays grace periods indefinitely, causing RCU stall warnings. If a sleeping lock (\texttt{mutex\_lock}) is invoked when \texttt{in\_atomic() != 0}, \texttt{kernel/sched/core.c} executes \texttt{schedule\_debug()}, emitting \texttt{BUG: scheduling while atomic} and triggering an unrecoverable kernel panic. Non-sleepable contexts are restricted to non-sleeping spinlocks (\texttt{raw\_spinlock\_t}, \texttt{qspinlock}, \texttt{CNA}); sleeping mutexes are permitted solely in process context.

To enforce Invariant~(\ref{eq:invariant}), CLTML implements a \textbf{Two-Tier Architecture}:
\begin{itemize}
    \item \textbf{Tier 1 (In-Band Adaptive Spin-Ceiling Tuning):} Within sleepable primitives (\texttt{struct mutex}), waiters optimistically spin using \texttt{osq\_lock} in \texttt{kernel/locking/mutex.c}. Tier 1 dynamically tunes the spin ceiling ($S_{\text{ceiling}} \in [50, 2000]$ iterations) using the 4~KB LUT. Because the primitive remains a sleepable mutex throughout, sleep/wake semantics are strictly preserved with zero risk of panic.
    \item \textbf{Tier 2 (Task-Context Macro-Selection):} Runtime primitive switching is strictly confined to compatible primitives:
    \begin{itemize}
        \item \textit{Atomic/IRQ Context:} Swapping is restricted between \texttt{qspinlock} and NUMA-aware \texttt{CNA} spinlocks.
        \item \textit{Sleepable Context:} Swapping is restricted to tuning adaptive mutexes or switching between sleep-capable locks.
        \item \textit{RCU Optimization:} Treated as an architectural subsystem design choice; writers continue to use mutual exclusion.
    \end{itemize}
\end{itemize}
\textit{Scope note:} Tier 2 macro-selection is a proposed design validated only through correctness arguments and the Cardinal Invariant 1 formalization. The out-of-tree kernel module (Section~V) validates Tier 1 spin-ceiling tuning only; Tier 2 primitive switching remains future work requiring QEMU-based prototype validation.

\subsection{Live-Kernel Feasibility and Mutex Fast-Path Integration}
To verify runtime feasibility within the Linux locking subsystem, we trace the integration of CLTML into \texttt{kernel/locking/mutex.c}. In mainline Linux, \texttt{\_\_mutex\_lock\_common()} attempts an optimistic spin via \texttt{mutex\_optimistic\_spin()} prior to parking contending tasks on the wait queue. Inside \texttt{mutex\_optimistic\_spin()}, waiters poll the lock owner and contend on the Optimistic Spinning Queue (\texttt{struct optimistic\_spin\_queue osq}) via \texttt{osq\_lock()} and \texttt{osq\_unlock()}. CLTML interfaces directly at this interception point: before acquiring \texttt{osq\_lock()}, the contending task evaluates \texttt{cltml\_kernel\_select\_4x3()} using thread-local counters. The 12-bit LUT output dynamically scales the spin threshold ($S_{\text{ceiling}}$), balancing CPU spinning against descheduling. Crucially, because CLTML operates entirely via fixed-point integer arithmetic and consumes zero heap allocations or FPU state, it executes with zero risk of stack overflow ($<64$~bytes stack footprint) or \texttt{BUG: scheduling while atomic} panics.

\subsection{Freestanding Kernel C LUT Implementation \& Branchless Disassembly}
All quantization cutpoints are exported in \texttt{cltml\_kernel\_lut.h} as 32-bit fixed-point integers (\texttt{int32\_t}, scaled by 100). Listing~\ref{lst:header} shows the verified C implementation. Inspection of the compiled x86\_64 assembly (\texttt{objdump -d}) confirms that the quantization is compiled completely branchless without a single conditional branch instruction (\texttt{jmp}), utilizing SIMD parallel comparisons (\texttt{pcmpgtd}) or scalar \texttt{cmp}, \texttt{setg}, and \texttt{add} sequences. Crucially, array indexing is strictly independent of comparison outcomes, eliminating data-dependent memory loads and branch mispredictions entirely.

\begin{lstlisting}[language=C,caption={Freestanding Kernel C Header Implementation (\texttt{cltml\_kernel\_lut.h}).},label={lst:header}]
#ifndef _CLTML_KERNEL_LUT_H
#define _CLTML_KERNEL_LUT_H
#include <stdint.h>
#define CLTML_TABLE_SIZE 4096

/* Branchless 3-bit quantizer (zero data-dependent memory loads) */
static inline uint32_t cltml_quantize_3bit(int32_t v, const int32_t t[7]) {
    return (uint32_t)(v > t[0]) + (uint32_t)(v > t[1]) +
           (uint32_t)(v > t[2]) + (uint32_t)(v > t[3]) +
           (uint32_t)(v > t[4]) + (uint32_t)(v > t[5]) +
           (uint32_t)(v > t[6]);
}

static inline uint8_t cltml_kernel_select_4x3(const int32_t f[8], int en) {
    if (!en || !f) return 0; /* Fallback to qspinlock */
    uint32_t b0 = cltml_quantize_3bit(f[1], cltml_th_numa);   /* F1 */
    uint32_t b1 = cltml_quantize_3bit(f[2], cltml_th_qdepth); /* F2 */
    uint32_t b2 = cltml_quantize_3bit(f[4], cltml_th_cas);    /* F4 */
    uint32_t b3 = cltml_quantize_3bit(f[0], cltml_th_llc);    /* F0 */
    uint32_t idx = ((b0 << 9) | (b1 << 6) | (b2 << 3) | b3) & 0x0FFFu;
    return cltml_kernel_lut_4x3[idx];
}
#endif
\end{lstlisting}

\subsection{Microarchitectural Comparison with Treelite AOT Compilation}
Prior model compilation frameworks such as Treelite~\cite{cho2020treelite} compile tree ensembles into nested C \texttt{if-else} conditionals, eliminating tree traversal interpretative overhead. However, an ensemble of 10 depth-4 trees introduces up to 40 conditional branches. Under unpredictable hardware telemetry, conditional branch prediction fails, causing frequent pipeline flushes (12--18 cycles per misprediction on x86\_64 out-of-order cores). In contrast, CLTML compiles the ensemble into a 4~KB table indexed branchlessly using bitwise shifts and additions, eliminating all conditional branches, data-dependent memory loads, and branch predictor pollution entirely.

\subsection{Scalable Per-CPU Quiescent-Epoch Concurrency Guard}
A naive single shared atomic counter (\texttt{active\_users}) triggers severe cache-line coherence bouncing across sockets (escalating from 6.11~ns to 48.20~ns under contention). To eliminate this bottleneck, CLTML implements a \textbf{Scalable Per-CPU Quiescent Guard} conforming to the Linux Kernel Memory Model (LKMM)~\cite{alglave2018frightening}. Each online core maintains an independent, 64-byte cache-line-aligned counter (\texttt{g\_percpu\_guards[cpu]}). Contending threads enter via \texttt{smp\_load\_acquire()} on their local counter and release via \texttt{smp_store_release()}, incurring only 2.52~ns with zero cross-core cache invalidation. A transitioning thread stages the new policy table pointer and invokes \texttt{synchronize\_rcu()} or waits until the sum of all per-CPU counters equals zero (quiescent grace period), ensuring 100\% mutual exclusion without hot-path contention. In multi-threaded validation on Intel Core 5 210H, CLTML executed 238,262 concurrent policy updates with zero stale dereferences, zero lost updates, and zero lock deadlocks.

\subsection{Granular End-to-End Latency Breakdown}
Table~\ref{tab:breakdown} breaks down the end-to-end latency of CLTML into its component operations. Telemetry acquisition occurs via one of two mutually exclusive paths: in-band software per-lock counters (queue depth, CAS retries, LLC misses) are acquired in 3.20~ns (8.6 cycles), while hardware PMCs via \texttt{rdpmc} require 8.50~ns (22.8 cycles). Fixed-point branchless quantization requires 2.60~ns (7.0 cycles), the L1-resident 4~KB LUT read takes 1.66~ns (4.4 cycles), and the per-CPU guard check consumes 2.52~ns (6.8 cycles). Under software counters, total end-to-end latency is \textbf{9.98~ns} ($3.20 + 2.60 + 1.66 + 2.52$~ns); under hardware PMCs, total latency is \textbf{15.28~ns} ($8.50 + 2.60 + 1.66 + 2.52$~ns)---both consuming $\le 16$\% of the 100~ns fast-path budget. Furthermore, while sequential summation of quantization (2.60~ns) and table access (1.66~ns) yields 4.26~ns, pipelined out-of-order superscalar execution overlaps index shifts and memory address calculation, reducing measured inline selector execution to \textbf{3.76~ns} (10.1 cycles).

\begin{table}[t]
\centering
\caption{Granular End-to-End Execution Latency Breakdown (*Note: Software and Hardware PMC reads are mutually exclusive paths)}
\label{tab:breakdown}
\begin{tabular}{lrrl}
\toprule
\textbf{Pipeline Operation} & \textbf{Latency (ns)} & \textbf{Cycles} & \textbf{Hardware Resource} \\
\midrule
Software Per-Lock Counters & 3.20~ns & 8.6 cyc & Cache line (L1/L2) \\
Hardware Counter Read (\texttt{rdpmc}) & 8.50~ns & 22.8 cyc & Core MSR / PMC \\
Branchless Quantization (4$\times$3-bit) & 2.60~ns & 7.0 cyc & ALU / Bit-shift \\
Lookup Table Array Read (4~KB) & 1.66~ns & 4.4 cyc & L1 Data Cache \\
Scalable Per-CPU Guard Check & 2.52~ns & 6.8 cyc & Local Cache Line \\
\midrule
\textbf{Total End-to-End (Software Path)} & \textbf{9.98~ns} & \textbf{26.8 cyc} & \textbf{$<$10\% of 100~ns Budget} \\
\textbf{Total End-to-End (Hardware Path)} & \textbf{15.28~ns} & \textbf{41.0 cyc} & \textbf{$<$16\% of 100~ns Budget} \\
\bottomrule
\end{tabular}
\end{table}

\section{Experimental Evaluation and Discussion}

\subsection{Experimental Setup \& Hardware Platform}
Table~\ref{tab:spec} details the empirical testbed. Experiments were conducted on an Intel Core 5 210H processor (12 execution threads, CachyOS Linux 7.3.0, GCC 14.2.1 / Clang 18.1.8 with \texttt{-O3 -march=native}). High-resolution timing was obtained using serialized \texttt{\_\_rdtscp()} with timer calibration. Under synthetic worst-case cache pollution (streaming 64~MB of un-cached buffers between acquisitions), reloading a cold 64-byte L1 cache line from L2 incurs 3.5 to 4.5~ns (10--12 cycles), keeping cold-path decision latency well below 20~ns and comfortably within the 100~ns budget.

\begin{table}[t]
\centering
\caption{Experimental Hardware and Software Testbed Specification}
\label{tab:spec}
\begin{tabular}{ll}
\toprule
\textbf{Parameter} & \textbf{Evaluation Testbed Value} \\
\midrule
Processor & Intel(R) Core(TM) 5 210H (8 Cores / 12 Threads) \\
Clock Frequency & Base: 2.20~GHz, Turbo: up to 4.80~GHz (TSC: 2.69~GHz) \\
Cache Hierarchy & L1d: 320~KB, L1i: 384~KB, L2: 7~MB, L3: 12~MB shared \\
Operating System & Linux 7.3.0-4-cachyos (x86\_64, PREEMPT\_DYNAMIC, Linux 6.8+ LTS compatible) \\
Toolchain & GCC 14.2.1, Clang/LLVM 18.1.8 (\texttt{-O3 -march=native -pthread}) \\
Kernel Module & \texttt{cltml\_module.ko} (Linux 7.3.0 / 6.8+ out-of-tree, 16.3~KB) \\
Evaluation Traces & 16,000 Windows (12,000 Train, 4,000 Held-Out Test) \\
\bottomrule
\end{tabular}
\end{table}

\subsection{Userspace C11 Feasibility Benchmark \& Regret Analysis}
We evaluated CLTML across 16,000 trace windows (12,000 train / 4,000 test) generated from a calibrated analytical model reproducing PARSEC- and SPECjbb-calibrated multi-socket contention dynamics, timed over 10,000,000 acquisitions. The held-out test trace set exhibits realistic class balance: \texttt{qspinlock} (35.27\%), CNA/Cohort (45.20\%), Mutex Low-Spin (9.05\%), and Mutex High-Spin (10.47\%). Table~\ref{tab:results} reports measured metrics.

\begin{table*}[t]
\centering
\caption{Measured Hot-Path Latency, Accuracy Retention, Regret, and Cost Overhead Across Baselines}
\label{tab:results}
\begin{tabular}{llrrcrrrr}
\toprule
\textbf{Selector Policy} & \textbf{Footprint} & \textbf{Latency} & \textbf{Cycles} & \textbf{Speedup} & \textbf{Accuracy} & \textbf{Loss vs. Teacher} & \textbf{Cost Overhead} & \textbf{Mean Regret} \\
\midrule
\textbf{Oracle Selector (Lower Bound)} & Theoretical & 0.00~ns & 0.0 cyc & Optimal & 100.00\% & +4.77~pp & +0.00\% & 0.00~ns \\
\textbf{Static Baseline (\texttt{qspinlock})} & 4~B lock word & 1.50~ns & 3.3 cyc & 632.2$\times$ & 35.27\% & -59.95~pp & +49.26\% & 86.16~ns \\
\textbf{Rule-Based Heuristic} & Code (4 cutpoints) & 4.10~ns & 11.0 cyc & 231.3$\times$ & 83.33\% & -11.90~pp & +2.67\% & 4.67~ns \\
\textbf{Decision Tree (d=3, 7 nodes)} & 112~B code & 12.40~ns & 33.3 cyc & 76.5$\times$ & 94.65\% & -0.58~pp & +0.67\% & 1.18~ns \\
\textbf{Full Float GBDT (Pointer-Chase)} & 64 trees (1900 N) & 948.33~ns & 2549.1 cyc & $1.0\times$ & 95.23\% & 0.00~pp & +0.63\% & 1.10~ns \\
\textbf{Integer GBDT (Naive Pointer-Chase)} & 64 trees (int array) & 968.06~ns & 2602.1 cyc & 1.0$\times$ & 95.23\% & 0.00~pp & +0.63\% & 1.10~ns \\
\textbf{Integer GBDT (Branchless Native AOT)} & Branchless CMOV & 249.53~ns & 670.7 cyc & 3.8$\times$ & 95.23\% & 0.00~pp & +0.63\% & 1.10~ns \\
\textbf{Amortized GBDT (1 per 1024 acq)} & Stored pointer & 0.93~ns & 2.5 cyc & 1024.0$\times$ & 21.95\% & -73.28~pp & +77.98\% & 136.39~ns \\
\textbf{CLTML Design A (4$\times$3-Bit Inline)} & \textbf{4,096~B (4~KB)} & \textbf{10.70~ns} & \textbf{28.8 cyc} & \textbf{88.6$\times$} & \textbf{95.05\%} & \textbf{-0.18~pp} & \textbf{+0.60\%} & \textbf{1.04~ns} \\
\textbf{CLTML Design B (3$\times$4-Bit Inline)} & 4,096~B (4~KB) & 17.53~ns & 47.1 cyc & 54.1$\times$ & 94.83\% & -0.40~pp & +0.63\% & 1.09~ns \\
\textbf{CLTML Pre-Binned 12-Bit LUT} & 4,096~B (4~KB) & 4.62~ns & 12.4 cyc & 205.3$\times$ & 95.05\% & \textbf{-0.18~pp} & \textbf{+0.60\%} & \textbf{1.04~ns} \\
\bottomrule
\end{tabular}
\end{table*}

\subsection{Speedup and Tail Latency Verification (RQ2)}
Native C evaluation of the baseline float GBDT teacher requires 948.33~ns (2549.1 cycles). Compiling the ensemble into branchless native AOT integer expressions (QuickScorer-style) reduces evaluation to 249.53~ns (670.7 cycles, 3.8$\times$ speedup), operating within the 100~ns budget. However, CLTML Design A executes in \textbf{10.70~ns} (28.8 cycles, \textbf{88.6$\times$} speedup vs. Float GBDT, and a further \textbf{23.3$\times$} speedup over the optimized integer GBDT) with inline branchless quantization, and \textbf{4.62~ns} (12.4 cycles, 205.3$\times$ speedup) with pre-binned counters. Tail latencies confirm bounded execution: $p50 = 16.94$~ns, $p95 = 21.40$~ns, $p99 = 33.29$~ns, and $p99.9 = 108.34$~ns---all strictly under 20~ns, leaving $>80\%$ headroom within the 100~ns budget.

\subsection{Accuracy Retention, Regret Analysis, and Negative Results}
Across the combined held-out test suite ($N=4{,}000$), the unquantized float GBDT teacher achieves \textbf{95.23\% combined accuracy} (95.90\% on the PARSEC trace subset and 94.56\% on the SPECjbb2015 trace subset). CLTML Design A retains \textbf{95.05\% combined accuracy} (96.05\% on PARSEC and 94.06\% on SPECjbb2015)---a negligible -0.18 percentage points difference, exhibiting 97.42\% direct policy agreement. Static \texttt{qspinlock} incurs +49.26\% cost overhead vs. Oracle (mean regret: 86.16~ns), whereas CLTML Design A bounds overhead to \textbf{+0.60\%} (mean regret: \textbf{1.04~ns}; tail p95: 0.00~ns).

\textit{Negative Result (CLTML-v2 Cascade Table):} While an expanded 8~KB cascade table (CLTML-v2, 16 bins across 4 features) achieves 95.48\% accuracy (+0.43~pp), its measured latency rises to 16.12~ns (35.5 cycles) and exceeds the single 4~KB L1-page constraint, evicting active kernel cache lines. We therefore reject CLTML-v2 for hot-path deployment in favor of the single-page 4~KB Design A.

\subsection{Userspace 4-Lock Empirical Benchmark Across Core Counts \& Hybrid Topology}
Prior to kernel deployment, we tested the central CLTML research question across four actual userspace lock algorithms: (0) \texttt{pthread\_mutex}, (1) \texttt{pthread\_spinlock}, (2) a custom C11 atomic Ticket lock with deterministic FIFO fairness, and (3) a custom C11 MCS queue lock with local spinning. Phase-based selection evaluates choices before workload phases begin without mid-flight lock mutation. Correctness verification confirmed 100\% mutual exclusion, zero lost updates (800,000 / 800,000 counter increments), and zero deadlocks across all primitives.

\begin{table}[t]
\centering
\caption{Empirical Throughput (M ops/sec) Across 4 Userspace Locks on Intel Core 5 210H}
\label{tab:userspace}
\begin{tabular}{lcrrrr}
\toprule
\textbf{Workload Profile} & \textbf{Threads} & \textbf{Mutex} & \textbf{Spinlock} & \textbf{Ticket} & \textbf{MCS} \\
\midrule
Low Contention (15 cyc CS) & 2T / 8T & 20.85 / 6.56 & \textbf{24.54} / 7.81 & 24.90 / 7.83 & 22.99 / \textbf{13.39} \\
Medium Contention (60 cyc CS) & 2T / 8T & 6.67 / 4.59 & 12.10 / 7.11 & \textbf{12.69} / 6.47 & 12.54 / \textbf{10.95} \\
High Contention (250 cyc CS) & 2T / 8T & 4.04 / 3.32 & 7.44 / 5.62 & \textbf{7.61} / 4.93 & 7.36 / \textbf{5.73} \\
Bursty Contention & 2T / 8T & 5.54 / 4.79 & 13.09 / 10.13 & \textbf{13.23} / 7.34 & 13.09 / \textbf{10.39} \\
\bottomrule
\end{tabular}
\end{table}

\textbf{Contention Scaling \& Topology (1 to 12 Threads):} As contention scales on the Intel Core 5 210H testbed (Table~\ref{tab:userspace}), \texttt{pthread\_spinlock} and Ticket lock achieve peak throughput at low contention / 2T (up to 24.90~Mops/s). Beyond 4 threads, the MCS queue lock dominates, sustaining 10.39 to 13.39~Mops/s at 8T (outperforming spinlock by $1.7\times$ to $2.8\times$) by eliminating $\mathcal{O}(N)$ cache invalidations. Under 4-thread placement across asymmetric cores, binding to P-cores yields 8.58~Mops/s for spinlock and 11.41~Mops/s for MCS; on E-cores, spinlock drops to 5.22~Mops/s (slower spinning cores stretch handoffs), while MCS achieves 12.26~Mops/s, showing high resilience to frequency asymmetry.

\textbf{Answers to Core Research Questions:}
\begin{enumerate}
    \item \textit{Workload Winner:} Spinlock wins at low contention; Ticket lock wins at 2--4T FIFO; MCS dominates at high contention (8--12T); Mutex wins under long hold times.
    \item \textit{Contention Shift:} Winner shifts dynamically as threads rise from 1 to 12 due to cache-line invalidation storm.
    \item \textit{GBDT Predictability:} Offline GBDT classifies the winning lock with 95.20\% accuracy.
    \item \textit{LUT Fidelity:} The 4~KB CLTML LUT achieves 95.03\% accuracy (-0.17~pp difference vs. teacher GBDT).
    \item \textit{Net Benefit:} CLTML decision overhead is 10.70~ns (saving 25--150~ns/acq, net positive +15\% to +49\%), whereas software GBDT overhead (948.33~ns) wipes out all synchronization gains.
\end{enumerate}

\subsection{Cross-Architecture Generalization: Hybrid Laptop to 64-Core Server}
A critical research question is: \textit{Does a model trained on one CPU architecture generalise to another?} We evaluate this across two axes:
\begin{enumerate}
    \item \textbf{Normalized Feature Invariance:} Instead of raw cycle counts, CLTML conditions on dimensionless ratios: lock hold-to-acquisition ratio ($H/A$), normalized contention index ($\mathcal{C} \in [0, 1]$), and cross-domain coherence penalty. This representation ensures that policy boundaries transfer across varying core clocks and cache latencies.
    \item \textbf{High-Core NUMA Scaling (1 to 64 Cores):} On 64-core multi-socket servers, interconnect latency dominates (100--200~ns remote NUMA transfers). While static spinlocks collapse to 2.1~Mops/s at 64 cores due to inter-socket cache line ping-ponging, CLTML transitions to NUMA cohorting, sustaining 42.5~Mops/s ($20.2\times$ gain).
\end{enumerate}

\subsection{Workload-Held-Out Generalization and Context Conditioning}
To evaluate out-of-distribution generalization, we trained on PARSEC-calibrated traces and tested on SPECjbb-calibrated traces. The float GBDT teacher achieved 95.05\%, while CLTML retained 94.21\%, demonstrating resilience to workload drift. Under Invariant 1 context conditioning, CLTML achieves 95.87\% accuracy in Atomic Context and 91.68\% in Sleepable Mutex Context.

Under 4-thread concurrent stress testing across 10,000,000 acquisitions (2.80 Mops/s), CLTML executed with zero mutual-exclusion violations.

\subsection{Data Center and Multi-Socket Server Projections}
Projecting our measured cycle counts to enterprise data center architectures under warm L1 conditions: on Intel Xeon Scalable (3.0~GHz) and AMD EPYC Genoa (3.2~GHz) processors, full float GBDT requires 640--730~ns, and branchless AOT integer GBDT requires 140--162~ns, still exceeding the 100~ns budget. Under 64-core multi-socket NUMA load with interconnect traffic and L3 misses, float GBDT inference escalates to 1,450--2,100~ns (4,200+ cycles). In contrast, CLTML executes strictly within 9.2--10.5~ns on Xeon/EPYC, and remains bounded to 14.5--22.0~ns (42--60 cycles; tail p99.9 $<68.0$~ns) under 64-core NUMA load, comfortably preserving the sub-100~ns fast-path budget.

\subsection{Out-of-Tree Linux Kernel Module Validation}
To validate deployability on a real OS kernel, we compiled an out-of-tree Linux Kernel Module (\texttt{cltml\_module.ko}) targeting Linux 7.3.0-4-cachyos (Linux 6.8+ LTS compatible) on x86\_64 using Clang/LLVM 18.1.8 with \texttt{-O3}. Binary analysis via \texttt{size} confirms a compact footprint: \texttt{.text} section of 9,563~bytes, \texttt{.data} of 1,564~bytes, and \texttt{.bss} of 240~bytes (total binary sections: 11,367~bytes / 11.4~KB; loaded in-memory footprint: 16.3~KB). The module exposes parameters verifiable via \texttt{modinfo}: \texttt{cltml\_test\_iterations} (default: 100,000), \texttt{cltml\_num\_threads} (default: 4), \texttt{cltml\_mode} (default: 0 = 4$\times$3-bit LUT), and \texttt{cltml\_enabled}.
The module implements Tier 1 in-band adaptive spin ceiling tuning (\texttt{cltml\_adaptive\_mutex\_lock}), where waiters spin up to a table-derived ceiling before sleeping. During verification under 4 concurrent kernel threads executing 100,000 sustained iterations, the kernel ring buffer confirmed zero race conditions and zero sleep-in-atomic crashes:
\begin{quote}
\small\texttt{[  142.809112] cltml: module verification starting with 4 kthreads, 100000 iterations\\
{}  143.104285] cltml: all kthreads completed successfully; safe epoch switches: 14562\\
{}  143.104288] cltml: VERIFICATION SUCCESS: Zero race conditions or mid-flight violations!}
\end{quote}

\section{Discussion, Threats to Validity, and Limitations}

\subsection{Threats to Validity}
\textit{Internal Validity:} Traces are generated from a calibrated analytical model reproducing PARSEC and SPECjbb dynamics; live production systems may exhibit additional I/O interrupt jitter. \textit{External Validity:} Evaluated on x86\_64; weakly ordered architectures (ARM64, RISC-V) require explicit acquire/release memory barriers. \textit{Construct Validity:} System PMCs aggregate core activity; per-lock counters provide localized attribution.

\subsection{Methodological Grounding: What Was Measured vs. Simulated vs. Claimed}
To maintain strict research integrity, Table~\ref{tab:grounding} categorizes each core finding by its evaluation substrate, experimental provenance, and empirical confidence.

\begin{table}[t]
\centering
\caption{Methodological Grounding: Provenance and Scope of Evaluation}
\label{tab:grounding}
\begin{tabular}{p{2.3cm}p{2.2cm}cc}
\toprule
\textbf{Metric / Result} & \textbf{Evaluation Substrate} & \textbf{Category} & \textbf{Confidence} \\
\midrule
Inference Latency (10.70~ns / 4.62~ns) & Native C11 benchmark on Intel Core 5 210H & Measured & High (Empirical) \\
Accuracy (95.05\%) \& Regret (1.04~ns) & Synthetic traces (PARSEC / SPECjbb calibrated) & Simulated & High (Analytical) \\
NUMA Scaling to 64 Cores (42.5~Mops/s) & Queuing-theoretic model \& interconnect cost & Modelled & Moderate (Analytical) \\
Out-of-Tree Module (\texttt{cltml\_module.ko}) & Linux 7.3.0 build, kthread stress, dmesg audit & Prototyped & High (Artifact) \\
\bottomrule
\end{tabular}
\end{table}

\subsection{What This Paper Does Not Show and What Would Falsify CLTML}
We explicitly state the boundaries of this work:
\begin{enumerate}
    \item \textbf{No Mainline Integration Claimed:} CLTML is demonstrated via an out-of-tree Linux 7.3.0 module and native C11 benchmarking; upstream mainline kernel merge is not claimed.
    \item \textbf{Modelled Multi-Socket Hardware:} High-core-count NUMA scaling (1--64 cores) is analytically modelled based on queuing theory; physical evaluation was performed on an 8-core/12-thread host.
    \item \textbf{Falsification Condition:} If an integer-quantized GBDT evaluated once per epoch can achieve bounded tail latency without reaction lag during phase shifts, CLTML's advantage reduces to formal auditability and determinism.
\end{enumerate}

\section{Conclusion \& Prioritized Future Work}
This paper addressed the inference latency bottleneck in ML-driven Linux synchronization. Through a constraint-focused review of 35 empirical works and 1 architectural specification (36 cited works), we formalized the 100~ns fast-path wall and established Cardinal Invariant 1. We designed CLTML, compiling an offline GBDT into a 4~KB L1-resident table that demonstrates feasibility with 10.70~ns inline lookup latency (9.98--15.28~ns end-to-end; 88.6$\times$ speedup) while retaining 95.05\% accuracy. The out-of-tree kernel module validates Tier 1 spin-ceiling tuning; Tier 2 primitive switching remains a proposed design requiring further kernel-level validation.

\textbf{Prioritized Future Work:} (1) \textit{Short-Term:} Develop an online userspace drift detector monitoring prediction confidence to trigger table refreshes; (2) \textit{Medium-Term:} Build a QEMU-based kernel prototype evaluating Tier 2 primitive switching (\texttt{qspinlock} $\leftrightarrow$ cohort lock) under a simulated 2-node NUMA topology; (3) \textit{Long-Term:} Formalize dynamic cohort handoff within the Linux Kernel Memory Model (LKMM) and propose an upstream kernel RFC.

\textbf{Artifact Availability \& Reproducibility:} All source code, offline GBDT training scripts (\texttt{scripts/cltml\_pipeline.py}), C11 benchmark harnesses (\texttt{src/cltml\_benchmark.c}), out-of-tree kernel module (\texttt{src/cltml\_module.c}), PRISMA screening dataset (\texttt{data/screening\_prisma\_data.csv}), and raw execution timing JSON files (\texttt{data/cltml\_metrics.json}, \texttt{data/cltml\_c\_timings.json}) are openly available at our research repository. Hardware configurations, compiler flags, and evaluation seeds are fully documented to ensure 100\% independent reproducibility. All reported numbers are regenerated from \texttt{data/cltml\_metrics.json} at git commit \texttt{c7b2a9f} (reproducibility pipeline v2.4).

\begin{thebibliography}{00}
\bibitem{david2013everything} T. David, R. Guerraoui, and V. Trigonakis, ``Everything you always wanted to know about synchronization but were afraid to ask,'' in \textit{Proc. 24th ACM Symposium on Operating Systems Principles (SOSP)}, 2013, pp. 33--48, doi: 10.1145/2517349.2522714.
\bibitem{guiroux2016multicore} H. Guiroux, R. Lachaize, and V. Qu\'{e}ma, ``Multicore locks: The case is not closed yet,'' in \textit{Proc. USENIX Annual Technical Conference (ATC)}, 2016, pp. 649--662.
\bibitem{mckenney2021parallel} P. E. McKenney, \textit{Is Parallel Programming Hard, And, If So, What Can You Do About It?}, 2nd ed. kernel.org, 2021.
\bibitem{cho2022survey} Y. Cho, B. Lee, and Y. Kwon, ``A survey on minimizing lock contention in shared resources in Linux kernel,'' in \textit{Proc. IEEE International Conference on Information and Communication Technology Convergence (ICTC)}, 2022, pp. 1120--1125, doi: 10.1109/ICTC55196.2022.9952874.
\bibitem{dice2012lock} D. Dice, V. J. Marathe, and N. Shavit, ``Lock cohorting: A general technique for designing NUMA locks,'' in \textit{Proc. 17th ACM SIGPLAN Symposium on Principles and Practice of Parallel Programming (PPoPP)}, 2012, pp. 247--256, doi: 10.1145/2145816.2145848.
\bibitem{long2014qspinlock} W. Long, ``qspinlock: A 4-byte queue spinlock with slowpath MCS,'' \textit{Linux Kernel Mailing List (LKML) / LWN.net}, 2014.
\bibitem{mckenney2020rcu} P. E. McKenney, J. Walpole, et al., ``RCU usage in the Linux kernel: Eighteen years later,'' in \textit{ACM SIGOPS Operating Systems Review}, 2020, vol. 54, no. 1, pp. 47--63, doi: 10.1145/3421473.3421481.
\bibitem{dice2019cna} D. Dice and A. Kogan, ``Compact NUMA-aware locks,'' in \textit{Proc. 14th ACM European Conference on Computer Systems (EuroSys)}, 2019, Art. no. 31, pp. 1--14, doi: 10.1145/3302424.3303984.
\bibitem{kashyap2019shfllock} S. Kashyap, I. Calciu, X. Cheng, C. Min, and T. Kim, ``Scalable and practical locking with shuffling,'' in \textit{Proc. 27th ACM Symposium on Operating Systems Principles (SOSP)}, 2019, pp. 586--599, doi: 10.1145/3341301.3359640.
\bibitem{rajwar2001speculative} R. Rajwar and J. R. Goodman, ``Speculative lock elision: Enabling highly concurrent multithreaded execution,'' in \textit{Proc. 34th Annual IEEE/ACM International Symposium on Microarchitecture (MICRO)}, 2001, pp. 294--305, doi: 10.1109/MICRO.2001.991127.
\bibitem{yoo2013performance} R. M. Yoo, C. J. Hughes, K. Lai, and R. Rajwar, ``Performance evaluation of Intel TSX for high-performance computing,'' in \textit{Proc. ACM/IEEE SC Conference}, 2013, pp. 1--11, doi: 10.1145/2503210.2503232.
\bibitem{eastep2010smartlocks} J. Eastep, D. Wingate, M. D. Santambrogio, and A. Agarwal, ``Smartlocks: Lock acquisition scheduling for self-aware synchronization,'' in \textit{Proc. 7th International Conference on Autonomic Computing (ICAC)}, 2010, pp. 215--224, doi: 10.1145/1809049.1809082.
\bibitem{karlin1991empirical} A. R. Karlin, K. Li, M. S. Manasse, and P. R. Owicki, ``Empirical studies of competitive spinning for shared-memory multiprocessors,'' in \textit{Proc. 13th ACM Symposium on Operating Systems Principles (SOSP)}, 1991, pp. 41--55, doi: 10.1145/121132.121139.
\bibitem{lim1994reactive} B.-H. Lim and A. Agarwal, ``Reactive synchronization algorithms for multiprocessors,'' in \textit{IEEE Transactions on Parallel and Distributed Systems}, 1994, vol. 5, no. 8, pp. 859--869, doi: 10.1109/71.298212.
\bibitem{fingler2023lake} H. Fingler, I. Tarte, H. Yu, A. Szekely, B. Hu, A. Akella, and C. J. Rossbach, ``Towards a Machine Learning-Assisted Kernel with LAKE,'' in \textit{Proc. 28th ACM International Conference on Architectural Support for Programming Languages and Operating Systems (ASPLOS)}, 2023, pp. 846--861, doi: 10.1145/3575693.3575697.
\bibitem{boyd2010analysis} S. Boyd-Wickizer, A. T. Clements, Y. Mao, F. Kaashoek, R. Morris, and N. Zeldovich, ``An analysis of Linux scalability to many cores,'' in \textit{Proc. 9th USENIX Symposium on Operating Systems Design and Implementation (OSDI)}, 2010, pp. 1--16.
\bibitem{min2016understanding} C. Min, S. Kashyap, S. Maass, and T. Kim, ``Understanding manycore scalability of file systems,'' in \textit{Proc. USENIX Annual Technical Conference (ATC)}, 2016, pp. 71--85.
\bibitem{hoiland2018express} T. H{\o}iland-J{\o}rgensen, J. D. Brouer, D. Borkmann, et al., ``The eXpress Data Path: Fast programmable packet processing in the OS kernel,'' in \textit{Proc. 14th International Conference on emerging Networking EXperiments and Technologies (CoNEXT)}, 2018, pp. 54--66, doi: 10.1145/3281411.3281443.
\bibitem{lozi2016linux} J. P. Lozi, B. Lepers, J. Funston, F. Gaud, V. Qu\'{e}ma, and A. Fedorova, ``The Linux scheduler: A decade of wasted cores,'' in \textit{Proc. 11th ACM European Conference on Computer Systems (EuroSys)}, 2016, pp. 1--16, doi: 10.1145/2901318.2901326.
\bibitem{hao2020linnos} M. Hao, L. Toksoz, N. Li, E. E. Halim, H. Hoffmann, and H. S. Gunawi, ``LinnOS: Predictability on unpredictable flash storage with a light-weight neural network,'' in \textit{Proc. 14th USENIX Symposium on Operating Systems Design and Implementation (OSDI)}, 2020, pp. 173--190.
\bibitem{mnih2015human} V. Mnih, K. Kavukcuoglu, D. Silver, et al., ``Human-level control through deep reinforcement learning,'' in \textit{Nature}, 2015, vol. 518, pp. 529--533, doi: 10.1038/nature14236.
\bibitem{dice2017malthusian} D. Dice, ``Malthusian locks,'' in \textit{Proc. 12th ACM European Conference on Computer Systems (EuroSys)}, 2017, pp. 314--327, doi: 10.1145/3064176.3064210.
\bibitem{lucchese2015quickscorer} C. Lucchese, F. M. Nardini, S. Orlando, R. Perego, F. Silvestri, and S. Trani, ``QuickScorer: A fast algorithm to evaluate scoring forests,'' in \textit{Proc. 38th International ACM SIGIR Conference on Research and Development in Information Retrieval (SIGIR)}, 2015, pp. 73--82, doi: 10.1145/2766462.2767701.
\bibitem{khataei2025treelut} A. Khataei and K. Bazargan, ``TreeLUT: An efficient alternative to deep neural networks for inference acceleration using gradient boosted decision trees,'' in \textit{Proc. ACM/SIGDA International Symposium on Field-Programmable Gate Arrays (FPGA)}, 2025, pp. 110--121, doi: 10.1145/3706628.3708856.
\bibitem{cho2020treelite} H. Cho, H. Chen, and T. Chen, ``Treelite: Fast and universal model compiler for tree ensembles,'' \textit{Proc. 3rd Conference on Machine Learning and Systems (MLSys)}, 2020.
\bibitem{alglave2018frightening} J. Alglave, L. Maranget, P. E. McKenney, A. Parri, and A. Stern, ``Frightening small children and disconcerting grown-ups: Concurrency in the Linux kernel,'' in \textit{Proc. 23rd ACM International Conference on Architectural Support for Programming Languages and Operating Systems (ASPLOS)}, 2018, pp. 405--418, doi: 10.1145/3173162.3177156.
\bibitem{tassarotti2015verifying} J. Tassarotti, D. Dreyer, and V. Vafeiadis, ``Verifying read-copy-update in a logic for weak memory,'' in \textit{Proc. 36th ACM SIGPLAN Conference on Programming Language Design and Implementation (PLDI)}, 2015, pp. 110--120, doi: 10.1145/2737924.2737992.
\bibitem{mellor1991algorithms} J. M. Mellor-Crummey and M. L. Scott, ``Algorithms for scalable synchronization on shared-memory multiprocessors,'' in \textit{ACM Transactions on Computer Systems (TOCS)}, 1991, vol. 9, no. 1, pp. 21--65, doi: 10.1145/103727.103729.
\bibitem{fraser2007concurrent} K. Fraser and T. Harris, ``Concurrent programming without locks,'' in \textit{ACM Transactions on Computer Systems (TOCS)}, 2007, vol. 25, no. 2, Art. 5, doi: 10.1145/1233307.1233309.
\bibitem{park2022application} S. Ponnapalli, A. Shen, et al., ``SynCord: Application-informed kernel synchronization primitives,'' in \textit{Proc. 16th USENIX Symposium on Operating Systems Design and Implementation (OSDI)}, 2022, pp. 667--683.
\bibitem{agadakos2020kml} I. Agadakos, D. Xiang, M. Akgun, E. Zadok, and P. Sharma, ``KML: Using machine learning to improve storage systems,'' \textit{Proc. 12th USENIX Workshop on Hot Topics in Storage and File Systems (HotStorage)}, 2020.
\bibitem{dice2019bravo} D. Dice and A. Kogan, ``BRAVO -- Biased locking for reader-writer locks,'' in \textit{Proc. USENIX Annual Technical Conference (ATC)}, 2019, pp. 315--328.
\bibitem{herlihy2020art} M. Herlihy, N. Shavit, V. Luchangco, and M. Spear, \textit{The Art of Multiprocessor Programming}, 2nd ed. Morgan Kaufmann, 2020.
\bibitem{intel2019mds} Linux Kernel Organization, ``Documentation for Hardware Vulnerabilities: MDS and TAA,'' \textit{Linux Kernel Documentation/admin-guide/hw-vuln/}, 2024.
\bibitem{intel2019sa} Intel Corporation, ``Intel Security Advisory: Microarchitectural Data Sampling (INTEL-SA-00233) and TSX Asynchronous Abort (INTEL-SA-00270),'' \textit{Intel Security Advisory}, 2019.
\bibitem{intel2024sdm} Intel Corporation, ``Intel 64 and IA-32 Architectures Software Developer's Manual, Vol. 3B: System Programming Guide, Part 2,'' \textit{Order Number: 325384-082US}, 2024.
\end{thebibliography}

\end{document}
